% =======================================================================
% edu-i-day 2026 (SVIA), ENTER Museum Derendingen, 25.11.2026
% Kurzvortrag FreeFlower, 20 Minuten inkl. Ausprobieren und Fragen
% =======================================================================
% Zahlen auf den Folien: Stand 22.11.2026 (git log, PDFs auf dem Server).
% Buchcover: Figures/extern/ ist nicht im Repo (Fremdbilder ohne freie
% Lizenz). Fehlen die Dateien, erscheint ein grauer Platzhalter.

\def\edupath{Various/eduiday_2026}

% Zeitangabe unten rechts, extrem klein: \zeit{Start}{Ende} am Anfang jedes Frames
\newcommand{\zeitangabe}{}
\newcommand{\zeit}[2]{\gdef\zeitangabe{#1--#2 / 20:00}}
\setbeamertemplate{footline}{%
	\hfill{\fontsize{4}{5}\selectfont\color{gray!70}\zeitangabe\hspace{3mm}\insertframenumber}\hspace*{2mm}%
	\vspace*{1.2mm}}

% Farben und Stile für die Grafiken
\colorlet{ffblue}{FDUblue}
\colorlet{ffgreen}{FDUgreen}
\colorlet{fforange}{Orange!85!black}
\colorlet{fflight}{ffblue!55}
\tikzset{
	box/.style={rounded corners=2pt, draw=#1, fill=#1!10, thick, align=center, font=\scriptsize, inner sep=3pt},
	box/.default=ffblue,
	pfeil/.style={-{Stealth[length=2mm]}, thick, gray!70!black},
	tile/.style={rounded corners=3pt, fill=#1!10, draw=#1!70, thick, minimum width=3.7cm, minimum height=1.7cm, text width=3.5cm, align=left, font=\scriptsize, inner sep=4pt},
	tile/.default=ffblue,
	thema/.style={rounded corners=2pt, fill=#1!12, draw=#1!60, font=\scriptsize, inner sep=2.5pt, minimum height=4.6mm, text depth=.4ex},
	gruppe/.style={font=\small\bfseries, text=#1},
	quelle/.style={rounded corners=2pt, fill=#1!12, draw=#1!60, align=center, inner sep=2.5pt, fill opacity=.92, text opacity=1},
	personbox/.style={box=#1, font=\scriptsize, text width=3.1cm, minimum height=1.2cm},
}

% Buchcover mit Platzhalter, falls die Datei lokal fehlt
\newcommand{\buchcover}[2]{%
	\IfFileExists{\edupath/Figures/extern/#1}{%
		\includegraphics[height=1.9cm]{\edupath/Figures/extern/#1}%
	}{%
		\tikz\node[draw=gray, fill=gray!10, minimum width=.235\linewidth, minimum height=1.8cm, text width=.2\linewidth, align=center, font=\tiny]{INFORMATIK\\[1mm]#2};%
	}%
}

% QR-Code mit Beschriftung
\newcommand{\qrkarte}[3]{%
	\begin{tikzpicture}
		\node[draw=ffblue!60, rounded corners=3pt, fill=white, inner sep=4pt, text width=2.55cm, minimum height=3.85cm, align=center] {%
			{\color{ffblue}\Large #1}\\[1mm]
			\qrcode[height=2.1cm]{#3}\\[1mm]
			{\scriptsize #2}};
	\end{tikzpicture}%
}

% Bild mit dünnem grauem Rahmen
\newcommand{\bild}[2]{\tikz[baseline=(b.base)]\node[draw=gray!50, fill=white, inner sep=0, line width=.4pt] (b) {\includegraphics[#1]{#2}};}

% Gemeinsame Höhe für Skriptseiten
\newcommand{\skriptseite}[2]{%
	\begin{tikzpicture}
		\node[inner sep=0, draw=gray!50, fill=white, drop shadow={opacity=.25}] (s) {\includegraphics[height=4.6cm]{\edupath/Figures/#1}};
		\node[below=1mm of s, font=\tiny, text width=2.6cm, align=center] {#2};
	\end{tikzpicture}%
}

\title{FreeFlower}
\subtitle{Informatiklehrmittel gemeinsam entwickeln nach dem Vorbild von Open Source}
\author{Cyril Blum}
\date{25.11.2026}

% -----------------------------------------------------------------------
% 1 Titel
% -----------------------------------------------------------------------
\begin{frame}[plain]
	\zeit{00:00}{00:30}
	\begin{tikzpicture}[remember picture, overlay]
		\fill[ffblue!6] (current page.north west) rectangle (current page.south east);
		\node[anchor=north east] at ([xshift=-8mm, yshift=-8mm]current page.north east) {\includegraphics[height=2.6cm]{Logos/FreeFlower.pdf}};
		\node[anchor=west, text width=10cm] at ([xshift=8mm, yshift=0mm]current page.west) {%
			{\color{ffblue}\fontsize{30}{32}\selectfont\bfseries FreeFlower}\\[3mm]
			{\large Informatiklehrmittel gemeinsam entwickeln\\nach dem Vorbild von Open Source}\\[7mm]
			{\normalsize Cyril Blum}\\[.5mm]
			{\footnotesize\color{gray!80!black} Kantonsschule im Lee Winterthur\\Kantonsschule Stadelhofen, Filiale Dübendorf}\\[5mm]
			{\footnotesize\color{gray!80!black} edu-i-day 2026, ENTER Museum Derendingen, 25.11.2026}};
		\node[anchor=south east, align=center, font=\tiny, text=gray!80!black] at ([xshift=-8mm, yshift=8mm]current page.south east) {%
			\qrcode[height=2.4cm]{https://github.com/CyrilBlum/FreeFlower}\\[1mm]
			github.com/CyrilBlum/FreeFlower};
	\end{tikzpicture}
\end{frame}

% -----------------------------------------------------------------------
% 2 Anekdote: Rückmeldungen zur Klett-Reihe der ETH
% -----------------------------------------------------------------------
\begin{frame}{Wie alles anfing: Rückmeldungen an den Verlag}
	\zeit{00:30}{01:45}
	\begin{columns}[T, onlytextwidth]
		\begin{column}{.49\textwidth}
			\centering
			\buchcover{klett_programmieren_robotik.jpg}{Programmieren und Robotik}\hfill%
			\buchcover{klett_data_science_sicherheit.jpg}{Data Science und Sicherheit}\hfill%
			\buchcover{klett_algorithmen_ki.jpg}{Algorithmen und Künstliche Intelligenz}\hfill%
			\buchcover{westermann_informatik.jpg}{Informatik für Schweizer Maturitätsschulen}\\[2mm]
			{\scriptsize Hromkovi\v{c} et al., \emph{Grundlagen der Informatik für Schweizer Maturitätsschulen}, Klett und Balmer\\
			\emph{Informatik für Schweizer Maturitätsschulen}, Westermann}\\[.5mm]
			{\tiny\color{gray} Cover: \copyright{} Klett und Balmer Verlag, Baar (Verlagsprospekt 2022), \copyright{} Westermann Verlag}
		\end{column}
		\begin{column}{.48\textwidth}
			\begin{tikzpicture}[node distance=2.5mm]
				\draw[gray!50, line width=1.2pt] (0,0) -- (0,-5.3);
				\foreach \y/\datum/\txt [count=\i from 2] in {
						0/18.09.2023/{18 Rückmeldungen zu zwei Bänden an Klett (Lösungen, Notation, Beispiele)},
						-1.45/21.09.2023/{Antwort: Korrektur \enquote{beim nächsten Nachdruck (das Datum ist offen)}},
						-2.9/19.10.2023/{Erster Commit im eigenen Repository (KLW)},
						-4.25/24.10.2023/{Mittagessen mit Juraj Hromkovi\v{c} und Regula Lacher}} {
						\fill[ffblue, visible on=<\i->] (0,\y) circle (1.3mm);
						\node[anchor=north west, font=\scriptsize\bfseries, text=ffblue, visible on=<\i->] at (.25,\y+.17) {\datum};
						\node[anchor=north west, font=\scriptsize, text width=5.1cm, visible on=<\i->] at (.25,\y-.2) {\txt};
					}
			\end{tikzpicture}
		\end{column}
	\end{columns}
	\vspace{1mm}
	\begin{center}
		\begin{tikzpicture}
			\node[box=fforange, font=\small, text width=11.5cm, visible on=<6->] {Gute Bücher. Korrigiert wird aber erst in der nächsten Auflage, und die Informatik ändert sich laufend.};
		\end{tikzpicture}
	\end{center}
\end{frame}

% -----------------------------------------------------------------------
% 2b Bewusst unübersichtliche Sammlung von Informatikquellen in der Schweiz
% Nur Namen als Text, keine Logos (Markenrechte). URLs am 25.11.2026 geprüft.
% -----------------------------------------------------------------------
\begin{frame}{Alle erfinden das Rad neu}
	\zeit{01:45}{02:15}
	\centering
	\resizebox{\textwidth}{!}{%
	\begin{tikzpicture}
		% Schritt/x/y/Drehung/Farbe/Schrift/Text, absichtlich durcheinander und teils überlappend
		\foreach \s/\x/\y/\r/\c/\f/\t in {
				1/0/0/-4/ffblue/\small/{Klett: Grundlagen der Informatik\\{\tiny ETH, 3 Bände}},
				1/5.2/.75/6/ffblue/\small/{Westermann: Informatik für\\Schweizer Maturitätsschulen},
				1/8.4/-.2/-8/ffblue/\footnotesize/{ABZ ETH\\{\tiny abz.inf.ethz.ch}},
				2/1.6/-1.3/9/ffgreen/\normalsize/{oinf.ch},
				2/5.4/-1.1/-3/ffgreen/\footnotesize/{informatikgarten.ch},
				2/9.6/-1.5/12/ffgreen/\footnotesize/{eduskript.org},
				2/-.6/-2.5/5/ffgreen/\footnotesize/{informatik.mygymer.ch},
				2/3.3/-2.7/-10/fforange/\small/{SVIA-Moodle},
				2/7.4/-2.6/4/ffgreen/\scriptsize/{Coding Labs\\{\tiny GitHub KMI-KPZ}},
				2/10.6/-3.1/-6/gray/\scriptsize/{inf-schule.de},
				2/.9/-3.9/-7/Purple/\footnotesize/{TigerJython},
				2/4.8/-3.8/8/Purple/\scriptsize/{jython.ch},
				2/8.3/-4.0/-2/fforange/\footnotesize/{Informatik-Biber},
				2/-.4/-5.1/10/gray/\scriptsize/{swisseduc.ch},
				2/2.9/-5.0/-5/fforange/\scriptsize/{iLearnIT.ch},
				2/6.1/-5.2/3/fforange/\scriptsize/{fit-in-IT},
				2/9.9/-5.0/9/gray/\scriptsize/{code.org},
				3/5.9/-.1/14/Red/\tiny/{Word-Dossier\\Kollegin A},
				3/10.9/.4/-12/Red/\tiny/{OneNote\\Kollege B},
				3/3.0/-.75/-14/Red/\tiny/{eigener Moodle-Kurs},
				3/7.0/-3.2/-15/Red/\tiny/{Skript als PDF\\von 2014},
				3/1.0/-4.6/18/Red/\tiny/{Jupyter-Notebooks\\Fachschaft C},
				3/11.0/-2.2/16/Red/\tiny/{Folien vom\\Vorgänger},
				4/4.1/-1.9/-18/ffblue/\tiny/{FreeFlower}} {
				\node[quelle=\c, font=\f, rotate=\r, visible on=<\s->] at (\x,\y) {\t};
			}
	\end{tikzpicture}}%

	\vspace{2mm}
	\visible<5->{\footnotesize Aus guten Gründen: kantonale Lehrpläne, eigene Schwerpunkte, Lizenzen ohne Anpassungsrecht.}
\end{frame}

% -----------------------------------------------------------------------
% 3 Einstieg: Früher war das Web offen (MySpace-Profil mit eigenem CSS)
% Anstoss: Brandon Thomas, How I cloned Adobe with AI, CODE Summit 2026
% Das Profil ist ein selbst gezeichnetes Mockup, kein Screenshot.
% Schritt 1 zeigt ein Standbild aus dem Talk (Figures/extern/, nicht im Repo).
% -----------------------------------------------------------------------
\begin{frame}[fragile]{Früher war vieles offen}
	\zeit{02:15}{03:00}
	\begin{columns}[c, onlytextwidth]
		\begin{column}{.5\textwidth}
			\centering
			\only<1>{%
				\IfFileExists{\edupath/Figures/extern/talk_brandon_thomas.jpg}{\bild{width=\linewidth}{\edupath/Figures/extern/talk_brandon_thomas.jpg}}{\tikz\node[draw=gray, fill=gray!10, minimum width=\linewidth, minimum height=3.6cm, font=\tiny]{Standbild Talk};}\\[1mm]
				{\tiny\color{gray} Standbild: Brandon Thomas, \emph{How I cloned Adobe with AI}, AI Engineer CODE Summit 2026 (YouTube)}}%
			\only<2->{%
			\begin{tikzpicture}[font=\tiny]
				% Browserfenster
				\draw[gray!60, fill=gray!15, rounded corners=2pt] (0,0) rectangle (6,4.6);
				\foreach \x/\c in {.2/red!60, .4/orange!70, .6/green!50!black} \fill[\c] (\x,4.42) circle (.06);
				\node[anchor=west, fill=white, draw=gray!40, rounded corners=1pt, inner sep=1.5pt, text width=4.8cm] at (.85,4.42) {\texttt{www.myspace.com/cyril}};
				% Profil mit eigenem CSS: Hintergrund, Farben, Schrift
				\fill[Purple!80!black] (.1,.1) rectangle (5.9,4.25);
				\foreach \x/\y in {.4/3.9, 1.3/1.0, 2.6/3.6, 3.9/.5, 5.3/2.9, 4.6/3.95, .8/2.1, 5.6/1.3, 2.2/.35} \node[text=yellow!80] at (\x,\y) {$\ast$};
				\node[draw=HotPink, line width=1pt, fill=black, minimum width=1.3cm, minimum height=1.3cm, text=white] at (1.05,3.15) {\faUser};
				\node[anchor=north west, text=HotPink, font=\scriptsize\bfseries] at (1.85,3.85) {Cyril};
				\node[anchor=north west, text=white, text width=2.6cm] at (1.85,3.45) {:-) \enquote{Online now!}\\Mood: kreativ};
				\node[anchor=north west, fill=black!70, text=Cyan, text width=2.3cm, inner sep=2pt] at (1.85,2.55) {\faMusic~Profilsong läuft \dots};
				\node[anchor=north west, text=yellow!80, font=\tiny\bfseries] at (.3,1.9) {Cyril's Top 8};
				\foreach \i in {0,...,3} \node[draw=HotPink, fill=black, text=white, minimum size=.55cm, inner sep=0] at (.6+\i*.7,1.25) {\faUser};
				% Feld \enquote{Über mich} mit eingefügtem CSS
				\node[anchor=north west, draw=gray!60, fill=white, text width=2.1cm, inner sep=2pt, font=\fontsize{4.5}{5.5}\selectfont\ttfamily, visible on=<3->] at (3.55,1.95) {%
					\textless{}style\textgreater\\
					body \{ background:\\
					~~url(sterne.gif); \}\\
					.name \{ color: hotpink; \}\\
					\textless{}/style\textgreater};
			\end{tikzpicture}\\[1mm]
			{\tiny\color{gray} Skizze eines MySpace-Profils um 2006, kein Screenshot}}
		\end{column}
		\begin{column}{.47\textwidth}
			\only<1>{{\small\bfseries \enquote{The Internet was once open and awesome}}\\[1mm]
				{\scriptsize Napster, GeoCities, Wikipedia, Flash: Alle konnten mitmachen und den Code anschauen.}\\[3mm]}
			\begin{itemize}\small
				\item<3-> Eigenes HTML und CSS im Feld \enquote{Über mich}
				\item<4-> Gefiel ein Profil, übernahm man den Code.
				\item<5-> \enquote{Quelltext anzeigen} zeigte jede Website.
			\end{itemize}
			\vspace{3mm}
			\visible<6->{\begin{tikzpicture}
				\node[box=fforange, font=\small, text width=5.4cm] {Das ginge auch bei Lehrmitteln.};
			\end{tikzpicture}}
			\vspace{3mm}

			\visible<1->{\tiny\color{gray} Anstoss: Brandon Thomas, \emph{How I cloned Adobe with AI}, AI Engineer CODE Summit 2026, \url{https://www.youtube.com/watch?v=UCL4K36SHp0}}
		\end{column}
	\end{columns}
\end{frame}

% -----------------------------------------------------------------------
% 3 Klassischer Weg vs. Open-Source-Weg
% -----------------------------------------------------------------------
\begin{frame}{Was passiert, wenn jemand einen Fehler findet?}
	\zeit{03:00}{04:15}
	\centering
	\begin{tikzpicture}[node distance=3mm, every node/.append style={minimum height=1.05cm, text width=1.95cm}]
		% klassischer Weg
		\node[gruppe=gray!70!black, anchor=west, text width=4cm, minimum height=0] at (-.2,1.15) {Gedrucktes Lehrmittel};
		\node[box=gray] (k1) at (1,0) {Fehler\\entdeckt};
		\node[box=gray, right=of k1, visible on=<2->] (k2) {Mail an\\Verlag};
		\node[box=gray, right=of k2, visible on=<3->] (k3) {Autorenteam\\prüft};
		\node[box=gray, right=of k3, visible on=<4->] (k4) {Nachdruck,\\neue Auflage};
		\node[box=gray, right=of k4, visible on=<5->] (k5) {Klassensatz\\neu kaufen};
		\foreach \a/\b [count=\i from 2] in {k1/k2,k2/k3,k3/k4,k4/k5} \draw[pfeil, visible on=<\i->] (\a) -- (\b);
		\node[font=\small\bfseries, text=gray!70!black, right=2mm of k5, text width=1.5cm, visible on=<6->] {Jahre};

		% FreeFlower
		\node[gruppe=ffgreen, anchor=west, text width=4cm, minimum height=0, visible on=<7->] at (-.2,-1.55) {FreeFlower};
		\node[box=ffgreen, below=1.6cm of k1, visible on=<7->] (f1) {SuS meldet Fehler im Forum};
		\node[box=ffgreen, right=of f1, visible on=<8->] (f2) {Korrektur in \texttt{.tex}};
		\node[box=ffgreen, right=of f2, visible on=<9->] (f3) {\texttt{git push}};
		\node[box=ffgreen, right=of f3, visible on=<10->] (f4) {Server kompiliert PDFs};
		\node[box=ffgreen, right=of f4, visible on=<11->] (f5) {Moodle zeigt neue Version};
		\foreach \a/\b [count=\i from 8] in {f1/f2,f2/f3,f3/f4,f4/f5} \draw[pfeil, visible on=<\i->] (\a) -- (\b);
		\node[font=\small\bfseries, text=ffgreen, right=2mm of f5, text width=1.5cm, visible on=<12->] {Minuten};
	\end{tikzpicture}

	\vspace{6mm}
	\small\visible<13->{Oft hat schon die nächste Klasse die korrigierte Version.}\\
	\visible<14->{Konstruktive Meldungen geben Punkte für die Zeugnisnote.}
\end{frame}

% -----------------------------------------------------------------------
% 4 Entstehung: Zeitstrahl und Commits
% -----------------------------------------------------------------------
\begin{frame}{Entstehung}
	\zeit{04:15}{05:30}
	\centering
	\begin{tikzpicture}
		\begin{axis}[
				ybar stacked, bar width=9mm,
				width=12.5cm, height=3.9cm,
				ymin=0, ymax=620,
				symbolic x coords={23-2,24-1,24-2,25-1,25-2,26-1,26-2},
				xticklabels={2.\,Hj.\ 2023, 1.\,Hj.\ 2024, 2.\,Hj.\ 2024, 1.\,Hj.\ 2025, 2.\,Hj.\ 2025, 1.\,Hj.\ 2026, 2.\,Hj.\ 2026},
				xtick=data, ticklabel style={font=\tiny}, ylabel={\scriptsize Commits}, ylabel style={yshift=-2mm},
				axis lines*=left, ymajorgrids, grid style={gray!20},
				legend style={font=\tiny, draw=none, at={(0.5,1.03)}, anchor=south, legend columns=2, /tikz/every even column/.append style={column sep=4mm}},
				enlarge x limits=0.08,
			]
			\addplot[fill=ffblue!70, draw=none] coordinates {(23-2,59) (24-1,326) (24-2,160) (25-1,543) (25-2,283) (26-1,385) (26-2,127)};
			\addlegendentry{ohne KI-Vermerk}
			\only<2->{%
				\addplot[fill=fforange, draw=none] coordinates {(23-2,0) (24-1,0) (24-2,0) (25-1,0) (25-2,0) (26-1,26) (26-2,58)};
				\addlegendentry{mit KI-Vermerk (Claude, Copilot)}}
		\end{axis}
	\end{tikzpicture}

	\vspace{1mm}
	\begin{tikzpicture}[x=1.62cm]
		\draw[gray!50, line width=1.2pt, -{Stealth}] (0,0) -- (6.6,0);
		\foreach \x/\datum/\txt [count=\i from 3] in {
				0/10.2023/{Start an der KLW\\(Name \enquote{LeeTeX})},
				1/06.2024/{PDFs vom eigenen Server},
				2/11.2024/{weitere Fächer, z.\,B. Geografie},
				3/03.2026/{Website mit Anleitungen},
				4/05.2026/{Name FreeFlower, öffentlich auf fit-in-IT},
				5/06.2026/{Artikel im SVIA-Interface},
				6/heute/{1967 Commits}} {
				\fill[ffblue, visible on=<\i->] (\x,0) circle (1.1mm);
				\node[above, font=\tiny\bfseries, text=ffblue, visible on=<\i->] at (\x,.08) {\datum};
				\node[below, font=\tiny, text width=1.5cm, align=center, visible on=<\i->] at (\x,-.1) {\txt};
			}
	\end{tikzpicture}
\end{frame}

% -----------------------------------------------------------------------
% 5 Wer arbeitet mit
% -----------------------------------------------------------------------
\begin{frame}{Wer arbeitet mit?}
	\zeit{05:30}{06:30}
	\centering
	\begin{tikzpicture}
		\node[inner sep=0] (logo) at (0,0) {\includegraphics[height=1.5cm]{Logos/FreeFlower.pdf}};
		\node[personbox=ffblue, visible on=<2->] (tg) at (-4.3,1.55) {\textbf{Thomas Graf}, KLW\\610 Commits\\Ergänzungsfach, Programmieren, Robotik};
		\node[personbox=ffblue, visible on=<3->] (cb) at (4.3,1.55) {\textbf{Cyril Blum}, KLW und FDU\\1306 Commits\\Grundlagenfach, Pipeline, Website};
		\node[personbox=fflight, visible on=<4->] (ma) at (-4.3,-.55) {\textbf{Michael Anderegg}, KLW\\23 Commits\\Robotik, Kryptologie};
		\node[personbox=fflight, visible on=<5->] (wei) at (4.3,-.55) {\textbf{weitere Beiträge}\\Beat Jäckle, Jonathan Rappl, JudyBee\\(2 bis 6 Commits)};
		\node[personbox=ffgreen, visible on=<6->] (sus) at (-2.2,-2.45) {\textbf{SchülerInnen}\\Meldungen im Moodle-Forum};
		\node[personbox=fforange, visible on=<7->] (ki) at (2.2,-2.45) {\textbf{KI-Agenten}\\84 Commits seit März 2026};
		\foreach \n [count=\i from 2] in {tg,cb,ma,wei,sus,ki} \draw[gray!50, thick, visible on=<\i->] (logo) -- (\n);
	\end{tikzpicture}

	\vspace{2mm}
	\visible<7->{\scriptsize\color{gray} \texttt{git log} bis 19.11.2026, Konten pro Person zusammengezählt}
\end{frame}

% -----------------------------------------------------------------------
% 6 Inhalte: Themenlandkarte
% -----------------------------------------------------------------------
\begin{frame}{Was ist schon drin?}
	\zeit{06:30}{08:00}
	\centering
	\resizebox{\textwidth}{!}{%
	\begin{tikzpicture}[node distance=1.2mm]
		% Grundlagenfach
		\node[gruppe=ffblue, anchor=west] (gf) at (0,0) {\faGraduationCap{} Grundlagenfach Informatik};
		\matrix[matrix of nodes, ampersand replacement=\&, anchor=north west, column sep=1.2mm, row sep=1.2mm, nodes={thema=ffblue, anchor=west}] at (0,-.35) {
			Programmieren (Python) \textbf{250} \& Computer und Kodierungen \textbf{77} \& Zahlendarstellungen \textbf{25} \\
			Datenintegrität \textbf{51} \& Kompression \textbf{19} \& Kryptologie \textbf{99} \\
			Netzwerke \textbf{80} \& Datenbanken (SQL) \textbf{65} \& Tabellenkalkulation \textbf{12} \\
			KI und Algorithmen \textbf{131} \& KI und Prompting \textbf{62} \& Randomisierte Algorithmen \textbf{26} \\
			Graphen (Folien) \& Game-Projekt (Pygame) \& Gesellschaft \\
		};
		% Ergänzungsfach
		\node[gruppe=ffgreen, anchor=west, visible on=<2->] at (12.7,0) {\faInfinity{} Ergänzungsfach};
		\matrix[matrix of nodes, ampersand replacement=\&, anchor=north west, row sep=1.2mm, nodes={thema=ffgreen, anchor=west, visible on=<2->}] at (12.7,-.35) {
			Endliche Automaten \textbf{50} \\
			Induktion und Rekursion \textbf{98} \\
			Kolmogorov-Komplexität \textbf{36} \\
			Komplexität (Folien) \\
			LaTeX (Folien) \\
		};
		% Freifach und andere
		\node[gruppe=fforange, anchor=west, visible on=<3->] at (0,-3.95) {\faRobot{} Freifach und Förderung};
		\matrix[matrix of nodes, ampersand replacement=\&, anchor=north west, column sep=1.2mm, nodes={thema=fforange, anchor=west, visible on=<3->}] at (0,-4.3) {
			Robotik \& Git-Tutorial \& Risiko simulieren \\
		};
		\node[gruppe=Purple, anchor=west, visible on=<4->] at (7.2,-3.95) {\faGlobeEurope{} andere Fächer};
		\matrix[matrix of nodes, ampersand replacement=\&, anchor=north west, column sep=1.2mm, nodes={thema=Purple, anchor=west, visible on=<4->}] at (7.2,-4.3) {
			Stadtgeografie \textbf{98} \& Geomorphologie \textbf{30} \& Statistik \textbf{27} \\
		};
	\end{tikzpicture}}%

	\vspace{3mm}
	\begin{tikzpicture}[zahl/.style={align=center, font=\tiny, text width=1.9cm}]
		\foreach \i/\z/\t [count=\k from 5] in {0/19/Skripte, 1/1238/Skriptseiten, 2/92/Foliensätze, 3/2537/Folien, 4/1078/Aufgaben, 5/627/TikZ-Grafiken} {
				\node[zahl, visible on=<\k->] at (\i*2.1,0) {{\color{ffblue}\Large\bfseries \z}\\\t};
			}
	\end{tikzpicture}\\
	\visible<10->{\tiny\color{gray} Fett: Seiten mit Lösungen, Stand 22.11.2026}
\end{frame}

% -----------------------------------------------------------------------
% 7 Ein Quelltext, viele Dokumente, automatische Veröffentlichung
% -----------------------------------------------------------------------
\begin{frame}{Ein Quelltext, viele Dokumente}
	\zeit{08:00}{09:00}
	\centering
	\begin{tikzpicture}[node distance=3.5mm and 9mm]
		\node[box, font=\small, text width=2.0cm, minimum height=1.6cm] (src) {\faFileCode\\[1mm]\texttt{Kapitel.tex}\\{\tiny Text, Aufgaben, Lösungen, TikZ}};
		\node[box=ffgreen, right=1.0cm of src, yshift=1.95cm, text width=3.0cm, visible on=<2->] (b1) {Skript mit Lösungen {\tiny(\texttt{book})}};
		\node[box=ffgreen, below=of b1, text width=3.0cm, visible on=<3->] (b2) {Skript ohne Lösungen {\tiny(\texttt{exerciseonly})}};
		\node[box=ffgreen, below=of b2, text width=3.0cm, visible on=<4->] (b3) {Folien {\tiny(\texttt{beamer})}};
		\node[box=ffgreen, below=of b3, text width=3.0cm, visible on=<5->] (b4) {Prüfungen, Lernkarten, Merkblätter};
		\foreach \b [count=\i from 2] in {b1,b2,b3,b4} \draw[pfeil, visible on=<\i->] (src.east) -- (\b.west);

		\node[box=gray, right=1.0cm of b2, yshift=-4.5mm, text width=2.5cm, minimum height=2.6cm, visible on=<6->] (pipe) {%
			\faGithub{} \texttt{git push}\\[1.2mm]
			\faCogs{} GitHub Actions\\[1.2mm]
			\faServer{} eigener Server\\{\tiny TeX Live 2024}\\[1.2mm]
			\faEnvelope{} Mail mit Änderungen};
		\foreach \b in {b1,b2,b3,b4} \draw[pfeil, gray!40, visible on=<6->] (\b.east) -- (pipe.west);

		\node[box=fforange, right=6mm of pipe, yshift=8mm, text width=2.3cm, visible on=<7->] (moodle) {\faLink{} feste PDF-Links in Moodle};
		\node[box=fforange, right=6mm of pipe, yshift=-8mm, text width=2.3cm, visible on=<8->] (web) {\faGlobe{} Website};
		\draw[pfeil, visible on=<7->] (pipe.east) -- (moodle.west);
		\draw[pfeil, visible on=<8->] (pipe.east) -- (web.west);
	\end{tikzpicture}

	\vspace{4mm}
	\small\visible<9->{Jede Änderung erneuert automatisch alle betroffenen PDFs.}
\end{frame}

% -----------------------------------------------------------------------
% 8 Blick in die Skripte
% -----------------------------------------------------------------------
\begin{frame}{Ein Blick in die Skripte}
	\zeit{09:00}{10:00}
	\centering
	% Programmieren: ganze Seiten 46 (TikZ-Grafik) und 47 (Aufgabe Wurzelspirale)
	\visible<1->{\skriptseite{skript_prog_funktion.pdf}{Programmieren: Funktion als Box (TikZ)}}\hfill
	\visible<1->{\skriptseite{skript_prog_aufgabe.pdf}{Programmieren: Aufgabe Wurzelspirale}}\hfill
	\visible<2->{\skriptseite{skript_addierer.pdf}{Computer und Kodierungen: 4-Bit-Addierer}}\hfill
	\visible<3->{\skriptseite{skript_hamming.pdf}{Datenintegrität: Überwachungstabelle}}

	\vspace{2mm}
	\visible<4->{\scriptsize Musterlösungen, farbige Boxen, Glossar und Lernziele im Anhang}
\end{frame}

% -----------------------------------------------------------------------
% 9 Textbasiert = bereit für Agenten
% -----------------------------------------------------------------------
\begin{frame}[fragile]{Reiner Text: lesbar für Menschen und Agenten}
	\zeit{10:00}{11:30}
	\begin{columns}[T, onlytextwidth]
		\begin{column}{.5\textwidth}
			{\scriptsize\bfseries \texttt{git diff} einer Forumsmeldung (30.10.2026)}
			\begin{tikzpicture}
				\node[draw=gray!40, fill=gray!5, rounded corners=2pt, inner sep=4pt, text width=6.1cm, font=\tiny\ttfamily, align=left] {%
					{\color{red!70!black}- ... ein Währungssystem von sieben verschiedene Münztypen.}\\
					{\color{green!45!black}+ ... ein Währungssystem von sieben verschiedenen Münztypen.}\\[1mm]
					{\color{red!70!black}- Wenn es sich um ein Schwarz-Weiss-Bild handelt ist dies ...}\\
					{\color{green!45!black}+ Wenn es sich um ein Schwarz-Weiss-Bild handelt, ist dies ...}};
			\end{tikzpicture}

			\vspace{3mm}
			\begin{tikzpicture}[node distance=3mm]
				\node[box, text width=1.6cm, visible on=<2->] (m) {\faUser\\Lehrperson beschreibt};
				\node[box=fforange, right=of m, text width=1.6cm, visible on=<3->] (a) {\faRobot\\Agent ändert \texttt{.tex}};
				\node[box, right=of a, text width=1.6cm, visible on=<4->] (p) {\faSearch\\prüft Diff und PDF};
				\draw[pfeil, visible on=<3->] (m) -- (a);
				\draw[pfeil, visible on=<4->] (a) -- (p);
				\draw[pfeil, visible on=<5->] (p.south) to[bend left=25] node[below, font=\tiny] {Korrektur} (a.south);
			\end{tikzpicture}
		\end{column}
		\begin{column}{.47\textwidth}
			\begin{tikzpicture}
				\node[tile=ffblue, text width=5.3cm, minimum height=0, visible on=<6->] {%
					\textbf{Von Hand geschrieben}\\
					Texte, Aufgaben und Lösungen seit 10.2023. 1883 von 1967 Commits ohne KI-Vermerk.};
			\end{tikzpicture}\\[2mm]
			\begin{tikzpicture}
				\node[tile=fforange, text width=5.3cm, minimum height=0, visible on=<7->] {%
					\textbf{Seit 2026 mit Agenten}\\
					TikZ-Grafiken, Folien zu bestehenden Kapiteln, Forumsmeldungen einarbeiten, Aufgabennummern abgleichen, abgeschnittenen Text finden.};
			\end{tikzpicture}\\[2mm]
			\begin{tikzpicture}
				\node[tile=ffgreen, text width=5.3cm, minimum height=0, visible on=<8->] {%
					\textbf{Was das möglich macht}\\
					Offener Quelltext, Textdateien, Git. Jede Änderung ist sichtbar und umkehrbar.};
			\end{tikzpicture}
		\end{column}
	\end{columns}
\end{frame}

% -----------------------------------------------------------------------
% 10 KI für TikZ: Beispiel Hamming-Folien
% -----------------------------------------------------------------------
% Makros aus Grundlagen_Info/05_Datenintegritaet/Slides/05_Hamming.tex übernommen
% Kartentrick-Codewort: a1..a16, dann c1..c9
\def\ktWerte{0,1,0,1,1,1,1,1,0,0,1,0,1,0,1,1,0,0,1,1,0,0,1,1,0}

% Überwachungstabelle der Kartentrick-Methode (16 Nachrichtenbits, 9 Kontrollbits)
\newcommand{\kartentrickUeberwachung}{%
    % Tabellenhintergrund
    \fill[blue!25] (0,0) rectangle (26,1);
    \fill[blue!15] (0,-9) rectangle (1,0);
    \fill[gray!8] (1,-9) rectangle (26,0);
    % Spalten- und Zeilenbeschriftung
    \foreach \i in {1,...,16} {
            \node at (\i+.5,.5) {\(a_{\i}\)};
        }
    \foreach \i in {1,...,9} {
            \pgfmathtruncatemacro{\col}{16+\i}
            \node at (\col+.5,.5) {\(c_{\i}\)};
            \node at (.5,.5-\i) {\(c_{\i}\)};
        }
    % c1 bis c4 überwachen die Spalten der Kartentrick-Tabelle
    \foreach \j in {1,...,4} {
            \foreach \r in {0,...,3} {
                    \pgfmathtruncatemacro{\col}{\j+4*\r}
                    \node at (\col+.5,.5-\j) {\(\times\)};
                }
        }
    % c5 bis c8 überwachen die Zeilen
    \foreach \r in {1,...,4} {
            \pgfmathtruncatemacro{\row}{\r+4}
            \foreach \j in {1,...,4} {
                    \pgfmathtruncatemacro{\col}{4*(\r-1)+\j}
                    \node at (\col+.5,.5-\row) {\(\times\)};
                }
        }
    % jedes Kontrollbit überwacht sich selbst
    \foreach \i in {1,...,8} {
            \pgfmathtruncatemacro{\col}{16+\i}
            \node at (\col+.5,.5-\i) {\(\times\)};
        }
    % c9 überwacht sämtliche Bits
    \foreach \col in {1,...,25} {
            \node at (\col+.5,-8.5) {\(\times\)};
        }
    \draw[step=1,black!30,thin] (0,-9) grid (26,1);
    \draw[black!65] (0,-9) rectangle (26,1);
}

% Zeile mit den Werten des Codeworts über der Tabelle; #1 = umgedrehtes Bit (0 = keines)
\newcommand{\kartentrickWertzeile}[1]{%
    \node[font=\scriptsize\itshape,anchor=east] at (1,1.5) {Wert};
    \foreach \col in {1,...,25} {
            \pgfmathtruncatemacro{\wert}{{\ktWerte}[\col-1]}
            \ifnum\col=#1
                \pgfmathtruncatemacro{\wert}{1-\wert}
                \node[magenta,font=\scriptsize\bfseries\ttfamily] at (\col+.5,1.5) {\wert};
            \else
                \node[font=\scriptsize\ttfamily] at (\col+.5,1.5) {\wert};
            \fi
        }
}

% TikZ-Schlüssel onslide aus 05_Hamming.tex: Optionen nur auf bestimmten Folien
\tikzset{onslide/.code args={<#1>#2}{\only<#1>{\pgfkeysalso{#2}}}}
% Frame aus 05_Hamming.tex übernommen (Übergang Kartentrick-Tabelle zu Überwachungstabelle)
\begin{frame}{Beispiel: von der Kartentrick-Tabelle zur Überwachungstabelle}
    \zeit{11:30}{12:00}
    \begin{center}
        \begin{tikzpicture}[every node/.style={font=\tiny}]
            % links: Kartentrick-Tabelle mit den Namen der Bits
            \begin{scope}[x=.5cm,y=.5cm]
                \foreach \r in {1,...,4} {
                        \foreach \c in {1,...,4} {
                                \pgfmathtruncatemacro{\nr}{4*(\r-1)+\c}
                                \fill[gray!12] (\c-1,-\r) rectangle (\c,-\r+1);
                                \node at (\c-.5,-\r+.5) {\(a_{\nr}\)};
                            }
                    }
                \foreach \c in {1,...,4} {
                        \fill[blue!20] (\c-1,0) rectangle (\c,1);
                        \node at (\c-.5,.5) {\(c_{\c}\)};
                    }
                \foreach \r in {1,...,4} {
                        \pgfmathtruncatemacro{\nr}{\r+4}
                        \fill[blue!20] (4,-\r) rectangle (5,-\r+1);
                        \node at (4.5,-\r+.5) {\(c_{\nr}\)};
                    }
                \fill[blue!20] (4,0) rectangle (5,1);
                \node at (4.5,.5) {\(c_9\)};
                \draw[step=1,black!65] (0,-4) grid (5,1);
                % a1 liegt in der Spalte von c1 und in der Zeile von c5
                \draw<2>[magenta,very thick] (0,-1) rectangle (1,0);
                \draw<2>[magenta,thick,dashed] (0,0) rectangle (1,1);
                \draw<2>[magenta,thick,dashed] (4,-1) rectangle (5,0);
                \coordinate (g-a1) at (0,-.5);
                % Bereich von c1: die erste Spalte
                \draw<3>[orange,very thick,rounded corners=2pt] (0,-4) rectangle (1,1);
                \coordinate (g-c1) at (1,.5);
            \end{scope}
            % rechts: Überwachungstabelle, erst ab Schritt 2 sichtbar (Platz bleibt reserviert, damit nichts springt)
            \begin{scope}[x=.4cm,y=.4cm,shift={(3cm,0)},onslide=<1>{opacity=0}]
                \kartentrickUeberwachung
                % Spalte von a1: Kreuze bei c1, c5, c9
                \draw<2>[magenta,very thick] (1,-9) rectangle (2,1);
                % dieselben Kontrollbits c1 und c5 wie links gestrichelt
                \draw<2>[magenta,thick,dashed] (1.15,-.85) rectangle (1.85,-.15);
                \draw<2>[magenta,thick,dashed] (1.15,-4.85) rectangle (1.85,-4.15);
                \coordinate (t-a1) at (1.5,-9);
                % Zeile von c1: a1, a5, a9, a13 und c1
                \draw<3>[orange,very thick] (0,-1) rectangle (26,0);
                \coordinate (t-c1) at (0,-.5);
            \end{scope}
            % Pfeil links um die Kartentrick-Tabelle herum zur Spalte von a1
            \draw<2>[->,magenta,very thick,rounded corners=4pt] (g-a1) -- (-.3cm,-.25cm) -- (-.3cm,-3.85cm) -| (t-a1);
            \draw<3>[->,orange,very thick] (g-c1) to[out=0,in=180] (t-c1);
        \end{tikzpicture}
    \end{center}
    \begin{itemize}
        \item<2-> \(a_1\) liegt in Spalte \(c_1\) und Zeile \(c_5\), \(c_9\) zählt alles $\rightarrow$ Spalte \(a_1\): \(c_1, c_5, c_9\)
        \item<3-> \(c_1\) prüft die erste Spalte $\rightarrow$ Zeile \(c_1\): \(a_1, a_5, a_9, a_{13}, c_1\)
    \end{itemize}
\end{frame}

\begin{frame}[fragile]{Beispiel: komplexe TikZ-Grafiken mit KI}
	\zeit{12:00}{12:30}
	\begin{columns}[T, onlytextwidth]
		\begin{column}{.4\textwidth}
\begin{lstlisting}[language={[LaTeX]TeX}, basicstyle=\tiny\ttfamily, frame=single, rulecolor=\color{gray!40}, backgroundcolor=\color{gray!5}, keywordstyle=\color{ffblue}, columns=fullflexible]
% Kartentrick-Tabelle mit Beispielwerten
\newcommand{\kartentrickGitter}[1]{%
  \foreach \r in {1,...,4} {
    \foreach \c in {1,...,4} {
      \pgfmathtruncatemacro{\nr}{4*(\r-1)+\c}
      \fill[gray!12] (\c-1,-\r)
        rectangle (\c,-\r+1);
      \pgfmathtruncatemacro{\wert}
        {{\ktWerte}[\nr-1]}
      \ifnum\nr=#1
        \pgfmathtruncatemacro{\wert}{1-\wert}
        \node[magenta] at (\c-.5,-\r+.5)
          {\texttt{\wert}};
      ...
\end{lstlisting}
			{\tiny\color{gray} \texttt{05\_Hamming.tex}, mit Claude erstellt (11.2026)}
		\end{column}
		\begin{column}{.58\textwidth}
			\centering
			% Überwachungstabelle mit denselben Schritten wie auf der Folie in 05_Hamming.tex
			\visible<2->{\resizebox{\linewidth}{!}{%
				\begin{tikzpicture}[x=.46cm,y=.4cm,every node/.style={font=\scriptsize}]
					\kartentrickUeberwachung
					\uncover<3->{\kartentrickWertzeile{0}}
					\uncover<4->{
						\node[font=\scriptsize\itshape] at (27.3,.5) {Einsen};
						\foreach \n [count=\i] in {2,2,4,4,2,4,2,4,14} {
								\node[green!50!black] at (27.3,.5-\i) {\n\ \checkmark};
							}
					}
					\draw<4->[orange,very thick] (0,-1) rectangle (26,0);
				\end{tikzpicture}}\\[1mm]
			{\tiny\color{gray} Folie \enquote{Überwachungstabelle}, animiert}}
		\end{column}
	\end{columns}
	\vspace{1mm}
	\visible<5->{\scriptsize Ich beschreibe, was die Grafik zeigen soll, und prüfe das Resultat. Von Hand bräuchte ich Stunden.}
\end{frame}

% -----------------------------------------------------------------------
% 10b Beispiel: animierte Folie Kasiski-Test
% Frame aus Grundlagen_Info/03_Kryptologie/Slides/L02d_Kryptoanalyse_Vigenere_Kasiski.tex
% -----------------------------------------------------------------------
\providecommand{\gold}[1]{\textcolor{Gold}{#1}}
\providecommand{\annotKasiski}[4]{\textcolor{#1}{\(\underset{\text{\tiny #4}}{\texttt{#2}}\)#3}}
\begin{frame}[fragile]{Beispiel: animierte Folie zum Kasiski-Test}{Kryptologie}
	\zeit{12:30}{13:00}
  \only<1->{\texttt{M\only<1-4>{EIN}\only<5->{\gold{EIN}}KL\only<1-4>{EIN}\only<5->{\red{EIN}}ERREIMSCH\only<1-4>{EIN}\only<5->{\red{EIN}}TK\only<1-4>{EIN}\only<5->{\green{EIN}}ERZUS\only<1-4>{EIN}\only<5->{\green{EIN}}}}

  % feste Höhe: Schlüssel, Alternative und durchgestrichene Alternative am selben Ort
  \begin{overprint}
    \onslide<1>\mbox{}
    \onslide<2-8,11->\texttt{C\only<1-4>{ODE}\only<5->{\gold{ODE}}CO\only<1-4>{DEC}\only<5->{\red{DEC}}ODECODECO\only<1-4>{DEC}\only<5->{\red{DEC}}OD\only<1-4>{ECO}\only<5->{\green{ECO}}DECOD\only<1-4>{ECO}\only<5->{\green{ECO}}}
    \onslide<9>\texttt{ABCDEFGHIJKLABCDEFGHIJKLABCDEFGHIJ??}
    \onslide<10>\sout{\texttt{ABCDEFGHIJKLABCDEFGHIJKLABCDEFGHIJ??}}
  \end{overprint}

  \visible<3->{\rule{0pt}{0pt}\(\vphantom{\underset{\text{\tiny 2}}{\texttt{S}}}\)\texttt{O\only<1-3>{SLR}\only<4->{\annotKasiski{Gold}{S}{LR}{2}}MZ\only<1-3>{HMP}\only<4->{\annotKasiski{red}{H}{MP}{7}}SUVGWPWEV\only<1-3>{HMP}\only<4->{\annotKasiski{red}{H}{MP}{19}}HN\only<1-3>{IKB}\only<4->{\annotKasiski{ForestGreen}{I}{KB}{24}}HVBIV\only<1-3>{IKB}\only<4->{\annotKasiski{ForestGreen}{I}{KB}{32}}}}\\[.5cm]

  \begin{itemize}
    \item<4-> \enquote{\red{HMP}} und \enquote{\green{IKB}} kommen je zweimal im Geheimtext vor

          \begin{itemize}
            \item<6-> \texttt{HMP}: Positionen 7 und 19. $\rightarrow$ 19-7 = 12 = $\blue{\mathbf{2\times 2}} \times 3$
                  \visible<12->{
                    \begin{tcolorbox}
                      $\mathbf{\underset{\text{Position 1 \texttt{HMP}}}{7}+\underset{\text{Codewort-Länge!}}{\blue{4}}\times \underset{\text{Anzahl Codewörter}}{3} = \underset{\text{Position 2 \texttt{HMP}}}{19}}$
                    \end{tcolorbox}
                  }
            \item<7-> \texttt{IKB}: Positionen 24 und 32. $\rightarrow$ 32-24 = 8  = $\blue{\mathbf{2\times 2}}\times 2$\visible<13->{
                    \begin{tcolorbox}
                      $\mathbf{\underset{\text{Position 1 \texttt{IKB}}}{24}+\underset{\text{Codewort-Länge!}}{\blue{4}}\times \underset{\text{Anzahl Codewörter}}{2} = \underset{\text{Position 2 \texttt{IKB}}}{32}}$
                    \end{tcolorbox}
                  }
            \item<8-> \textbf{Könnte die Schlüssellänge auch 12 sein?} \visible<10->{Nein, denn dann würden die beiden \texttt{IKB} nicht auf den selben Schlüsselteil fallen!}
            \item<11-> Distanz zwischen Positionen = Vielfaches der Schlüssellänge?
          \end{itemize}
  \end{itemize}
\end{frame}

% -----------------------------------------------------------------------
% 12 Ideen für den eigenen Unterricht
% -----------------------------------------------------------------------
\begin{frame}{Was Sie davon übernehmen könnten}
	\zeit{13:00}{13:45}
	\centering
	\begin{tikzpicture}[x=4.1cm, y=2.55cm, every node/.append style={minimum height=2.3cm}]
		\node[tile=ffblue, visible on=<1->] at (0,0) {{\color{ffblue}\Large\faComments}\\[.5mm]\textbf{Fehler-Forum in Moodle}\\SuS melden Fehler, gute Meldungen geben Punkte.};
		\node[tile=ffblue, visible on=<2->] at (1,0) {{\color{ffblue}\Large\faToggleOn}\\[.5mm]\textbf{Lösungen per Schalter}\\Eine Quelle, Skript mit und ohne Lösungen.};
		\node[tile=ffblue, visible on=<3->] at (2,0) {{\color{ffblue}\Large\faCodeBranch}\\[.5mm]\textbf{Git als Unterrichtsthema}\\Echte Commits zeigen, wie Software entsteht.};
		\node[tile=ffgreen, visible on=<4->] at (.5,-1) {{\color{ffgreen}\Large\faCheckDouble}\\[.5mm]\textbf{Automatische Kontrollen}\\Aufgabennummern, abgeschnittener Text, Kompilierfehler.};
		\node[tile=ffgreen, visible on=<5->] at (1.5,-1) {{\color{ffgreen}\Large\faDatabase}\\[.5mm]\textbf{Moodle-Fragen als Dateien}\\Fragesammlungen versioniert und überprüfbar.};
	\end{tikzpicture}
\end{frame}

% -----------------------------------------------------------------------
% 13 Grenzen
% -----------------------------------------------------------------------
\begin{frame}{Was (noch) nicht gut läuft}
	\zeit{13:45}{14:30}
	\centering
	\begin{tikzpicture}[x=6.1cm, y=2.0cm, every node/.append style={minimum height=1.75cm}]
		\node[tile=fforange, text width=5.4cm, minimum width=5.6cm, visible on=<1->] at (0,0) {\textbf{LaTeX ist eine Hürde}\\Zum Unterrichten reichen die PDFs. Mitschreiben braucht LaTeX und Git (oder einen Agenten).};
		\node[tile=fforange, text width=5.4cm, minimum width=5.6cm, visible on=<2->] at (1,0) {\textbf{KI-Fehler klingen plausibel}\\Jede Agenten-Änderung muss jemand mit Fachwissen prüfen.};
		\node[tile=fforange, text width=5.4cm, minimum width=5.6cm, visible on=<3->] at (0,-1) {\textbf{Wenige tragen die Wartung}\\Über 95\,\% der Commits stammen von zwei Personen.};
		\node[tile=fforange, text width=5.4cm, minimum width=5.6cm, visible on=<4->] at (1,-1) {\textbf{Bildrechte}\\Fremdbilder nur mit Lizenz und Quelle, das kostet Zeit.};
	\end{tikzpicture}
\end{frame}

% -----------------------------------------------------------------------
% 14 Ausprobieren
% -----------------------------------------------------------------------
\begin{frame}{Jetzt ausprobieren}
	\zeit{14:30}{17:15}
	\centering
	\qrkarte{\faFilePdf}{PDFs online lesen}{https://cyrilblum.github.io/FreeFlower/compiled-files.html}\hfill
	\qrkarte{\faCode}{Quelltext im Browser (VS Code)}{https://github.dev/CyrilBlum/FreeFlower}\hfill
	\qrkarte{\faDownload}{Repository als ZIP}{https://github.com/CyrilBlum/FreeFlower/archive/refs/heads/main.zip}

	\vspace{4mm}
	% Auftrag als schlichte Box: Kennzeichnung oben, Text darunter linksbündig
	\only<2->{%
		\begin{tcolorbox}[center, width=11.6cm, colback=ffblue!6, colframe=ffblue!55, boxrule=.6pt, arc=2pt, left=3mm, right=3mm, top=2mm, bottom=2mm, fontupper=\small]
			\aufgabebox{Zu zweit, etwa 4 Minuten}\par\smallskip
			Blättern Sie in einem Thema, das Sie unterrichten.\par\smallskip
			Was würden Sie übernehmen, was fehlt, was stört?
		\end{tcolorbox}}

	\vspace{1mm}
	\visible<2->{\tiny\color{gray} Auf GitHub öffnet die Taste \texttt{.} denselben Editor.}
\end{frame}

% -----------------------------------------------------------------------
% 15 Offene Diskussion
% -----------------------------------------------------------------------
\begin{frame}{Offene Fragen}
	\zeit{17:15}{20:00}
	\begin{columns}[c, onlytextwidth]
		\begin{column}{.34\textwidth}
			\centering
			\includegraphics[height=2.6cm]{Logos/FreeFlower.pdf}
		\end{column}
		\begin{column}{.62\textwidth}
			\begin{tikzpicture}[node distance=2.5mm]
				\node[box, text width=6.4cm, align=left, font=\small, visible on=<2->] (q1) {Was würden Sie übernehmen, was nicht, und weshalb?};
				\node[box, text width=6.4cm, align=left, font=\small, below=of q1, visible on=<3->] (q2) {Was müsste passieren, damit nicht alle das Rad neu erfinden?};
				\node[box, text width=6.4cm, align=left, font=\small, below=of q2, visible on=<4->] (q3) {Was bräuchten Sie, um mitzuschreiben?};
				\node[box=fforange, text width=6.4cm, align=left, font=\small, below=of q3, visible on=<5->] (q4) {Wo ziehen Sie die Grenze für KI-Agenten bei Lehrmitteln?};
			\end{tikzpicture}

			\vspace{5mm}
			\visible<6->{{\scriptsize Später gerne per Mail: cyril.blum@edu.zh.ch\\ cyrilblum.github.io/FreeFlower}}
		\end{column}
	\end{columns}
\end{frame}

% -----------------------------------------------------------------------
% Reserve: Mitmachen
% -----------------------------------------------------------------------
\begin{frame}{Reserve: So können Sie mitmachen}
	\gdef\zeitangabe{Reserve}
	\centering
	\begin{tikzpicture}[node distance=5mm, every node/.append style={text width=2.3cm, minimum height=1.3cm}]
		\node[box] (a) {\faEye\\PDFs nutzen};
		\node[box, right=of a, visible on=<2->] (b) {\faCodeBranch\\Fork oder Branch anlegen};
		\node[box, right=of b, visible on=<3->] (c) {\faEdit\\Anpassung in \texttt{.tex}};
		\node[box=ffgreen, right=of c, visible on=<4->] (d) {\faShare\\Pull Request};
		\foreach \x/\y [count=\i from 2] in {a/b,b/c,c/d} \draw[pfeil, visible on=<\i->] (\x) -- (\y);
	\end{tikzpicture}

	\vspace{6mm}
	\visible<5->{\small Lizenz: Unterlagen CC BY-NC-SA 4.0, Code und LaTeX-Klassen GPL-3.0-or-later.\\
		Regeln für die Zusammenarbeit: \url{https://cyrilblum.github.io/FreeFlower/zusammenarbeit.html}}
\end{frame}
